% block: preconditioning-heading pages:163,164 type:section
\subsection{{预条件技术}}\label{sec:preconditioning-techniques}
{迭代方法的收敛速度强烈依赖于系数矩阵的谱性质, 即特征值的分布. 因此, 自然可以尝试将原线性方程组变换为具有相同解而谱性质更有利的方程组. 若 $K$ 为非奇异矩阵, 则变换后的线性方程组}
\begin{equation}\label{eq:left-preconditioned-system}K^{-1}A\bm x=K^{-1}\bm b\end{equation}
{与原方程组具有相同的解. 因此, 将 Krylov 子空间方法应用于变换后的线性方程组, 就可得到原线性方程组的解. 矩阵 $K$ 称为预条件子, 良好的预条件子满足下列性质:}
\begin{itemize}\item $K^{-1}A\approx I$.
\item {对任意向量 $\bm z$, 容易得到 $K^{-1}\bm z$.}\end{itemize}
% block: split-preconditioning page:164 type:ordinary-paragraph
{若 $A$ 为 Hermitian 正定 (h.p.d.) 矩阵, 则变换后的方程组 \eqref{eq:left-preconditioned-system} 在实践中并不适用, 因为系数矩阵 $K^{-1}A$ 不再为 h.p.d. 矩阵. 为保留系数矩阵的 h.p.d. 结构, 将预条件子分解为 $K=K_1K_1^{\mathrm H}$, 再将线性方程组变换为}
\[K_1^{-1}AK_1^{-\mathrm H}(K_1^{\mathrm H}\bm x)=K_1^{-1}\bm b.\]
{上述系数矩阵 $K_1^{-1}AK_1^{-\mathrm H}$ 为 h.p.d. 矩阵. 因此, 可将用于 h.p.d. 线性方程组的 Krylov 子空间方法, 例如 CG 方法, 应用于该方程组. 若系数矩阵 $A$ 为非 Hermitian 矩阵, 则相应的预条件形式为}
\begin{equation}\label{eq:split-preconditioned-system}K_1^{-1}AK_2^{-1}(K_2\bm x)=K_1^{-1}\bm b,\end{equation}
{其中 $K=K_1K_2$. 上述形式称为左右预条件. 当 $K_1=I$ 时, 有}
\begin{equation}\label{eq:right-preconditioned-system}AK^{-1}(K\bm x)=\bm b.\end{equation}
{这种预条件称为右预条件; 当 $K_2=I$ 时, 得到 \eqref{eq:left-preconditioned-system}, 称为左预条件.}
% block: preconditioning-residual-choice page:164 type:ordinary-paragraph
{GMRES 方法和 GCR 方法经常采用右预条件. 在这种情况下, 这些方法寻找使原线性方程组 $A\bm x=\bm b$ 的残差的 2-范数极小的近似解. 另一方面, 若采用左预条件, 则这些方法寻找使变换后的线性方程组 $K^{-1}A\bm x=K^{-1}\bm b$ 的残差的 2-范数极小的近似解.}
{关于预条件技术的综述, 例如见 \BookCite{20}, \BookCite{146}. 另见双 Lanczos 型 Krylov 子空间方法的预条件算法 \BookCite{109}. 文献 \BookCite{49}, \BookCite{94}, \BookCite{136}, \BookCite{206} 研究了 Arnoldi 型 Krylov 子空间方法的预条件算法, 即用于奇异线性方程组的 GMRES 方法.}
% block: incomplete-decomposition-heading page:164 type:subsection
\subsubsection{{不完全矩阵分解预条件子}}\label{subsec:incomplete-decomposition-preconditioners}
{在 20 世纪的最后几十年中, 人们提出了许多预条件子. 在各种预条件子中, 最著名的一类基于系数矩阵的不完全分解, 可表示为 $A\approx K=LU$ (其中 $L$ 和 $U$ 为非奇异三角矩阵).}
% block: dilu-definition page:165 type:ordinary-paragraph
{最简单的不完全分解之一是 Pommerell 提出的 D-ILU 预条件子 (例如见 \BookCite{147}, pp.77--78). 该方法的思想如下: 首先, 将系数矩阵分解为对角部分, 严格下三角部分和严格上三角部分, 即 $A=D_A+L_A+U_A$; 其次, 以 $K=(D+L_A)D^{-1}(D+U_A)$ 为预条件子, 其中 $D$ 由 $\operatorname{diag}(K)=D_A$ 确定, 即对所有 $i$ 有 $(K)_{ii}=(D_A)_{ii}$. 因此, 只需计算 $D$.}
{D-ILU 预条件子的算法见 Algorithm\ \ref{alg:dilu}, 其中 $N_0(A)$ 为矩阵 $A$ 的非零元素的指标集, 即}
\[N_0(A):=\{(i,j):a_{i,j}\ne0\}.\]
{运行 Algorithm\ \ref{alg:dilu} 后, 得到 $d_1,d_2,\ldots,d_N$. 再令}
\[D=\begin{bmatrix}d_1^{-1}&&\\&\ddots&\\&&d_N^{-1}\end{bmatrix}\]
{就得到形式为 $K=(D+L_A)D^{-1}(D+U_A)$ 的 D-ILU 预条件子.}
% block: dilu-algorithm page:165 type:ordinary-paragraph
\begin{algorithm}[H]\caption{{D-ILU 分解}}\label{alg:dilu}\begin{algorithmic}[1]
\Statex\textbf{Input:} $A\in\mathbb C^{N\times N},\ N_0(A)$ {($A$ 的非零元素指标集)}
\Statex\textbf{Output:} $D^{-1}=\operatorname{diag}(d_1,d_2,\ldots,d_N)$
\For{$i=1,2,\ldots,N$}
\State $d_i=a_{i,i}$
\EndFor
\For{$i=1,2,\ldots,N$}
\State $d_i=1/d_i$
\For{$j=i+1,i+2,\ldots,N$}
\State {if $(i,j)\in N_0(A)$ 且 $(j,i)\in N_0(A)$, then $d_j=d_j-a_{j,i}d_i a_{i,j}$}
\EndFor\EndFor\end{algorithmic}\end{algorithm}
% block: dilu-implementation page:165 type:ordinary-paragraph
{下面说明预条件 Krylov 子空间方法的实现. D-ILU 预条件 Krylov 子空间方法需要计算 $K^{-1}$ 与向量 $\bm v$ 的乘积, 即 $\bm w=K^{-1}\bm v=(D+U_A)^{-1}D(D+L_A)^{-1}\bm v$, 向量 $\bm w$ 按下列步骤计算:}
\begin{enumerate}
\item {以前代求解 $(D+L_A)\bm y=\bm v$;}
\item {计算 $\bm z=D\bm y$;}
\item {以回代求解 $(D+U_A)\bm w=\bm z$.}
\end{enumerate}
{注意, 由于 D-ILU 预条件子显式包含原矩阵的非对角部分, 可以使用 Eisenstat 技巧 \BookCite{53} (另见 Section\ \ref{subsec:stationary-preconditioners}) 更高效地实现 D-ILU 预条件 Krylov 子空间方法.}
% block: ilu-zero-introduction page:166 type:ordinary-paragraph
{Meijerink 和 van der Vorst \BookCite{127} 给出了最著名且最有效的不完全 LU 分解之一. 该分解称为 ILU(0), 下面加以说明. ILU(0) 预条件子基于 Algorithm\ \ref{alg:lu-decomposition} 中的 $A$ 的 LU 分解. 为便于介绍 ILU(0) 预条件子, Algorithm\ \ref{alg:lu-row-oriented} 对 Algorithm\ \ref{alg:lu-decomposition} 略作了修改.}
% block: lu-in-place-storage page:166 type:ordinary-paragraph
{Algorithm\ \ref{alg:lu-row-oriented} 的输入为下列 $A$.}
\[A=\begin{bmatrix}a_{1,1}&a_{1,2}&\cdots&A_{1,N}\\a_{2,1}&a_{2,2}&\cdots&A_{2,N}\\\vdots&\vdots&\ddots&\vdots\\a_{N,1}&a_{N,2}&\cdots&a_{N,N}\end{bmatrix},\]
{输出为 $a_{i,j}$, 其中}
\[A=LU,\]
{这里}
\[L=\begin{bmatrix}1&&&\\a_{2,1}&\ddots&&\\\vdots&\ddots&\ddots&\\a_{N,1}&\cdots&a_{N,N-1}&1\end{bmatrix},\quad U=\begin{bmatrix}a_{1,1}&a_{1,2}&\cdots&a_{1,N}\\&\ddots&\ddots&\vdots\\&&\ddots&a_{N-1,N}\\&&&a_{N,N}\end{bmatrix}.\]
{注意, 输入的 $a_{i,j}$ 为原矩阵, 输出时 $a_{i,j}$ 被覆盖, 用来存储矩阵 $L$ 和 $U$.}
% block: row-oriented-lu-algorithm page:166 type:ordinary-paragraph
\begin{algorithm}[H]\caption{{LU 分解 (Algorithm\ \ref{alg:lu-decomposition} 的略微修改版本)}}\label{alg:lu-row-oriented}\begin{algorithmic}[1]
\For{$i=2,3,\ldots,N$}
\For{$k=1,2,\ldots,i-1$}
\State $a_{i,k}=a_{i,k}/a_{k,k}$
\For{$j=k+1,k+2,\ldots,N$}
\State $a_{i,j}=a_{i,j}-a_{i,k}\times a_{k,j}$
\EndFor\EndFor\EndFor\end{algorithmic}\end{algorithm}
% block: ilu-zero-definition pages:166,167 type:ordinary-paragraph
{ILU(0) 预条件子就是 Algorithm\ \ref{alg:lu-row-oriented} 中的 LU 分解, 但若 $(i,j)\notin N_0(A)$, 就不计算 $(i,j)$ 元素. ILU(0) 预条件子的算法见 Algorithm\ \ref{alg:ilu-zero}. 由 Algorithm\ \ref{alg:ilu-zero}, 分解后矩阵的非零元素数通常与原矩阵相同. 因此, 当 $A$ 为稀疏矩阵时, ILU(0) 预条件子的上三角矩阵和下三角矩阵也是稀疏矩阵. 这一事实对存储开销和计算开销都很重要.}
% block: ilu-zero-algorithm page:167 type:ordinary-paragraph
\begin{algorithm}[H]\caption{{ILU(0) 预条件子}}\label{alg:ilu-zero}\begin{algorithmic}[1]
\For{$i=2,3,\ldots,N$}
\For{$k=1,2,\ldots,i-1$}
\State if $(i,k)\in N_0(A)$, then $a_{i,k}=a_{i,k}/a_{k,k}$
\For{$j=k+1,k+2,\ldots,N$}
\State if $(i,j)\in N_0(A)$, then $a_{i,j}=a_{i,j}-a_{i,k}\times a_{k,j}$
\EndFor\EndFor\EndFor\end{algorithmic}\end{algorithm}
% block: fill-level-definition page:167 type:ordinary-paragraph
{作为 ILU(0) 预条件子的推广, 同一篇论文 \BookCite{127} 还提出了 ILU($p$) 预条件子, 其中 $p$ 表示填充级别. 填充级别的定义如下: 当 $a_{i,j}\ne0$ 时, 令位置 $(i,j)$ 的初始填充级别为零; 当 $a_{i,j}=0$ 时, 令位置 $(i,j)$ 的初始填充级别为 $\infty$. 则位置 $(i,j)$ 的填充级别定义为}
\begin{equation}\label{eq:ilu-fill-level}\operatorname{Level}_{ij}:=\min_{1\le k\le\min\{i,j\}}\{\operatorname{Level}_{ik}+\operatorname{Level}_{kj}+1\}.\end{equation}
% block: fill-level-example-original page:167 type:ordinary-paragraph
{例如, 考虑}
\begin{equation}\label{eq:ilu-fill-example-matrix}A=\begin{bmatrix}
*&*&0&0&0&*&0&0&0\\
*&*&*&0&0&0&*&0&0\\
0&*&*&*&0&0&0&*&0\\
0&0&*&*&*&0&0&0&*\\
0&0&0&*&*&*&0&0&0\\
*&0&0&0&*&*&*&0&0\\
0&*&0&0&0&*&*&*&0\\
0&0&*&0&0&0&*&*&*\\
0&0&0&*&0&0&0&*&*
\end{bmatrix},\end{equation}
{其中符号 ``*'' 表示非零元素. 由于 $A$ 的非零元素的级别为零, 而 $A$ 中零元素的级别为 $\infty$, 得到下列以填充级别为元素的矩阵:}
% block: fill-level-example-initial page:168 type:equation
\begin{equation}\label{eq:ilu-fill-initial-levels}\begin{bmatrix}
0&0&\infty&\infty&\infty&0&\infty&\infty&\infty\\
0&0&0&\infty&\infty&\infty&0&\infty&\infty\\
\infty&0&0&0&\infty&\infty&\infty&0&\infty\\
\infty&\infty&0&0&0&\infty&\infty&\infty&0\\
\infty&\infty&\infty&0&0&0&\infty&\infty&\infty\\
0&\infty&\infty&\infty&0&0&0&\infty&\infty\\
\infty&0&\infty&\infty&\infty&0&0&0&\infty\\
\infty&\infty&0&\infty&\infty&\infty&0&0&0\\
\infty&\infty&\infty&0&\infty&\infty&\infty&0&0
\end{bmatrix},\end{equation}
% block: fill-level-example-one page:168 type:ordinary-paragraph
{将 \eqref{eq:ilu-fill-level} 应用于 \eqref{eq:ilu-fill-initial-levels}, 得到下列位置 $(i,j)$ 的级别为 1 的矩阵:}
\[\begin{bmatrix}
0&0&\infty&\infty&\infty&0&\infty&\infty&\infty\\
0&0&0&\infty&\infty&1&0&\infty&\infty\\
\infty&0&0&0&\infty&\infty&1&0&\infty\\
\infty&\infty&0&0&0&\infty&\infty&1&0\\
\infty&\infty&\infty&0&0&0&\infty&\infty&1\\
0&1&\infty&\infty&0&0&0&\infty&\infty\\
\infty&0&1&\infty&\infty&0&0&0&\infty\\
\infty&\infty&0&1&\infty&\infty&0&0&0\\
\infty&\infty&\infty&0&1&\infty&\infty&0&0
\end{bmatrix},\]
% block: fill-level-example-two page:168 type:ordinary-paragraph
{将 \eqref{eq:ilu-fill-level} 应用于上述矩阵, 就可确定 $\operatorname{Level}_{ij}=2$ 的 $(i,j)$ 元素.}
\[\begin{bmatrix}
0&0&\infty&\infty&\infty&0&\infty&\infty&\infty\\
0&0&0&\infty&\infty&1&0&\infty&\infty\\
\infty&0&0&0&\infty&2&1&0&\infty\\
\infty&\infty&0&0&0&\infty&2&1&0\\
\infty&\infty&\infty&0&0&0&\infty&2&1\\
0&1&2&\infty&0&0&0&\infty&2\\
\infty&0&1&2&\infty&0&0&0&\infty\\
\infty&\infty&0&1&2&\infty&0&0&0\\
\infty&\infty&\infty&0&1&2&\infty&0&0
\end{bmatrix},\]
% block: fill-level-example-final pages:168,169 type:ordinary-paragraph
{将 \eqref{eq:ilu-fill-level} 应用于上述矩阵, 就确定了所有填充级别, 相应矩阵为}
\begin{equation}\label{eq:ilu-fill-final-levels}\begin{bmatrix}
0&0&\infty&\infty&\infty&0&\infty&\infty&\infty\\
0&0&0&\infty&\infty&1&0&\infty&\infty\\
\infty&0&0&0&\infty&2&1&0&\infty\\
\infty&\infty&0&0&0&3&2&1&0\\
\infty&\infty&\infty&0&0&0&3&2&1\\
0&1&2&3&0&0&0&3&2\\
\infty&0&1&2&3&0&0&0&3\\
\infty&\infty&0&1&2&3&0&0&0\\
\infty&\infty&\infty&0&1&2&3&0&0
\end{bmatrix},\end{equation}
% block: ilu-level-index-sets page:169 type:ordinary-paragraph
{设 $N_p(A)$ 为对应于 $\operatorname{Level}_{ij}\le p$ 的指标集, 即}
\[N_p(A):=\{(i,j):\operatorname{Level}_{ij}\le p\}.\]
{例如, 由 \eqref{eq:ilu-fill-final-levels}, 当 $p=1,2,3$ 时, \eqref{eq:ilu-fill-example-matrix} 的 $N_p(A)$ 如下:}
\[\begin{aligned}
N_1(A)&=N_0(A)\cup\{(2,6),(3,7),(4,8),(5,9),(6,2),(7,3),(8,4),(9,5)\},\\
N_2(A)&=N_1(A)\cup\{(3,6),(4,7),(5,8),(6,9),(6,3),(7,4),(8,5),(9,6)\},\\
N_3(A)&=N_2(A)\cup\{(4,6),(5,7),(6,8),(7,9),(6,4),(7,5),(8,6),(9,7)\}.
\end{aligned}\]
{ILU($p$) 预条件子就是 Algorithm\ \ref{alg:lu-row-oriented} 中的 LU 分解, 但若 $(i,j)\notin N_p(A)$, 就不计算 $(i,j)$ 元素. ILU($p$) 预条件子的算法见 Algorithm\ \ref{alg:ilu-level}.}
% block: ilu-level-algorithm page:169 type:ordinary-paragraph
\begin{algorithm}[H]\caption{{ILU($p$) 预条件子}}\label{alg:ilu-level}\begin{algorithmic}[1]
\For{$i=2,3,\ldots,N$}
\For{$k=1,2,\ldots,i-1$}
\State if $(i,k)\in N_p(A)$, then $a_{i,k}=a_{i,k}/a_{k,k}$
\For{$j=k+1,k+2,\ldots,N$}
\State if $(i,j)\in N_p(A)$, then $a_{i,j}=a_{i,j}-a_{i,k}\times a_{k,j}$
\EndFor\EndFor\EndFor\end{algorithmic}\end{algorithm}
% block: ilut-threshold page:169 type:ordinary-paragraph
{ILU 的另一种成功变型是 Saad \BookCite{150} 提出的采用阈值策略的 ILU. 该分解称为 ILUT($p,\tau$), 其中 $p$ 用于节省存储空间, $\tau$ 为舍弃元素的准则, 即将模较小的元素强制置零.}
% block: approximate-inverse-heading page:170 type:subsection
\subsubsection{{近似逆预条件子}}\label{subsec:approximate-inverse-preconditioners}
{前述预条件子的性能取决于分解所得矩阵的乘积 $\widetilde L\widetilde U$ 与 $A$ 的接近程度. 这里介绍判断良好预条件子的另一准则, 即预条件矩阵与 $A^{-1}$ 的接近程度. 基于这一准则, Grote 和 Huckle \BookCite{82} 尝试极小化下列 Frobenius 范数:}
\[\min_K\|I-AK\|_{\mathrm F},\]
{其中 $\|A\|_{\mathrm F}=\sqrt{\sum_{i,j}a_{i,j}^2}$. 由于上述极小化问题可写为}
\[\min_K\left\|[\bm e_1-A\bm k_1,\bm e_2-A\bm k_2,\ldots,\bm e_N-A\bm k_N]\right\|_{\mathrm F},\]
{就得到 $N$ 个相互独立的最小二乘问题}
\[\min_{\bm k_i}\|\bm e_i-A\bm k_i\|,\quad i=1,\ldots,N.\]
{因此, 该预条件子的构造非常适合并行计算.}
% block: ainv-biorthogonal-factorization page:170 type:ordinary-paragraph
{Benzi 和 Tůma \BookCite{24} 给出了关于近似逆的另一个出色思想, 即寻找非奇异矩阵 $V$ 和 $W$, 使得}
\begin{equation}\label{eq:ainv-biorthogonal-factorization}W^{\mathrm H}AV=D.\end{equation}
{由 $(W^{\mathrm H}AV)^{-1}=V^{-1}A^{-1}W^{-\mathrm H}=D^{-1}$ 得, 矩阵 $A$ 的逆为}
\[A^{-1}=VD^{-1}W^{\mathrm H}.\]
{由于等式 \eqref{eq:ainv-biorthogonal-factorization} 意味着当 $i\ne j$ 时 $(\bm w_i,A\bm v_j)=0$, 因而通过 $A$-双正交化过程计算 $V$ 和 $W$. 该算法称为稀疏近似逆 (AINV) 预条件子, 见 Algorithm\ \ref{alg:ainv}. 在算法中, $\bm a_i$ 和 $\bm c_i$ 分别表示 $A$ 和 $A^{\mathrm H}$ 的第 $i$ 列. 由 Algorithm\ \ref{alg:ainv} 可见, 该算法采用 MGS 风格的 $A$-双正交化过程, 用于非 Hermitian 矩阵. 对于 Hermitian 矩阵和正规矩阵, 也已发展出相应的近似逆预条件子, 见 \BookCite{21}, \BookCite{22}.}
% block: polynomial-preconditioner-heading page:171 type:subsection
\subsubsection{{矩阵多项式预条件子}}\label{subsec:polynomial-preconditioners}
{已经介绍了两类预条件子. 一类近似系数矩阵, 即 $A\approx M$, 且 $M^{-1}\bm v$ 易于计算. 另一类近似系数矩阵的逆, 即 $A^{-1}\approx M$. 这里介绍与稀疏近似逆密切相关的另一种方法. 由于该方法基于矩阵多项式, 因而称为多项式预条件子.}
% block: ainv-algorithm page:171 type:ordinary-paragraph
\begin{algorithm}[H]\caption{{AINV 预条件子}}\label{alg:ainv}\begin{algorithmic}[1]
\Statex\textbf{Input:} $\bm w_i^{(0)}=\bm v_i^{(0)}=\bm e_i,\ 1\le i\le N$
\For{$i=1,\ldots,N$}
\For{$j=i,\ldots,N$}
\State $p_j^{(i-1)}=\bm a_i^{\mathrm H}\bm v_j^{(i-1)},\quad q_j^{(i-1)}=\bm c_i^{\mathrm H}\bm w_j^{(i-1)}$
\EndFor
\State if $i=N$, then {退出}.
\For{$j=i+1,\ldots,N$}
\State $\bm v_j^{(i)}=\bm v_j^{(i-1)}-\left(\dfrac{p_j^{(i-1)}}{p_i^{(i-1)}}\right)\bm v_i^{(i-1)},\quad\bm w_j^{(i)}=\bm w_j^{(i-1)}-\left(\dfrac{q_j^{(i-1)}}{q_i^{(i-1)}}\right)\bm w_i^{(i-1)}$
\State $v_{k,j}^{(i)}=0$ if $|v_{k,j}^{(i)}|<\mathrm{Tol},\ 1\le k\le N$,\quad $w_{k,j}^{(i)}=0$ if $|w_{k,j}^{(i)}|<\mathrm{Tol},\ 1\le k\le N$
\EndFor\EndFor
\State $\bm z_i=\bm z_i^{(i-1)},\ \bm w_i=\bm w_i^{(i-1)},\ p_i=p_i^{(i-1)},\ 1\le i\le N$
\State $V=[\bm v_1,\bm v_2,\ldots,\bm v_N],\ V=[\bm w_1,\bm w_2,\ldots,\bm w_N],\ D=\operatorname{diag}(p_1,p_2,\ldots,p_N)$
\end{algorithmic}\end{algorithm}
% block: neumann-euler-preconditioners page:171 type:ordinary-paragraph
{最著名的多项式预条件子包括 Neumann 展开和 Euler 展开. 设 $A$ 具有形式 $I-B$, 且 $\rho(B)<1$. 则可将 $A$ 的逆矩阵写为}
\begin{center}\begin{tabular}{r@{}ll}
$A^{-1}$&$=\displaystyle\sum_{i=0}^{\infty}B^i$&{(Neumann 展开)}\\[8pt]
&$=\displaystyle\prod_{i=0}^{\infty}(I+B^{2^i}).$&{(Euler 展开)}
\end{tabular}\end{center}
{因此, 由上述展开, 利用 Neumann (或 Euler) 展开的低阶项, 就可容易地得到并进而控制 $A$ 的近似逆. 文献 \BookCite{50} 研究了这一方法, 另见 \BookCite{116}, \BookCite{198}.}
% block: stationary-preconditioner-heading page:172 type:subsection
\subsubsection{{基于定常迭代方法的预条件子}}\label{subsec:stationary-preconditioners}
{Section\ \ref{sec:stationary-iterations} 中定常迭代方法的思想也可用于为 Krylov 子空间方法构造高效预条件子. 本节介绍 Jacobi, Gauss--Seidel 和 SOR 型预条件子.}
% block: jacobi-preconditioner page:172 type:ordinary-paragraph
{Jacobi 预条件子 (或对角预条件子) 是最简单的预条件子, 通常比用不带预条件的 Krylov 子空间方法求解原线性方程组 $A\bm x=\bm b$ 更有效. 利用系数矩阵的对角元素, 按下式构造 Jacobi 预条件子:}
\[K_J:=\operatorname{diag}(a_{1,1},a_{2,2},\ldots,a_{n,n}).\]
{若矩阵 $A$ 为对角矩阵, 则 $K_J^{-1}A=I$. 这意味着 Jacobi 预条件 Krylov 子空间方法只需一步迭代就可得到解. 一般而言, 有一种方法是选取对角矩阵 $D$ 为预条件子, 使 $D^{-1}A$ 的条件数尽可能小. 关于 $D^{-1}A$ 的条件数分析见 \BookCite{59}, \BookCite{194}.}
% block: sgs-preconditioner page:172 type:ordinary-paragraph
{对称 Gauss--Seidel (SGS) 预条件子比 Jacobi 预条件子使用更多关于 $A$ 的信息. SGS 预条件子以 Gauss--Seidel 方法命名, 后者是 Section\ \ref{sec:stationary-iterations} 中的定常迭代方法之一. 该预条件子定义为}
\[K_{SGS}:=(D-L)D^{-1}(D-U),\]
{其中 $A:=D-L-U$. 若将 $A$ 的所有对角元素缩放为一, 就得到更简单的预条件子:}
\[K_{SGS}=(I-L)(I-U).\]
{在这种情况下, 已知有如下高效实现方式: 使用预条件子 $K_{SGS}$ 时, 得到下列形式的变换后线性方程组}
\[(I-L)^{-1}A(I-U)^{-1}\widetilde{\bm x}=\widetilde{\bm b},\quad\widetilde{\bm x}=(I-U)\bm x,\quad\widetilde{\bm b}=(I-L)^{-1}\bm b.\]
{因此, 将 Krylov 子空间方法应用于变换后的线性方程组时, 需要计算下列矩阵向量乘积:}
\[(I-L)^{-1}A(I-U)^{-1}\bm z.\]
% block: eisenstat-sgs-derivation pages:172,173 type:ordinary-paragraph
{当矩阵向量乘积 $A\bm z$ 的开销在每步迭代中占主导地位时, 这通常会使每次迭代的计算开销大约加倍. 然而, 回顾 $A=I-L-U$, 得}
\[\begin{aligned}(I-L)^{-1}A(I-U)^{-1}&=(I-L)^{-1}(I-L-U)(I-U)^{-1}\\
&=(I-L)^{-1}[(I-L)+(I-U-I)](I-U)^{-1}\\
&=(I-L)^{-1}[(I-L)(I-U)^{-1}+I-(I-U)^{-1}]\\
&=(I-U)^{-1}+(I-L)^{-1}[I-(I-U)^{-1}].\end{aligned}\]
{于是得到}
\[\begin{aligned}(I-L)^{-1}A(I-U)^{-1}\bm z&=(I-U)^{-1}\bm z+(I-L)^{-1}[I-(I-U)^{-1}]\bm z\\&=\bm t+(I-L)^{-1}(\bm z-\bm t),\end{aligned}\]
{其中 $\bm t:=(I-U)^{-1}\bm z$. 因此可见, 上述运算的开销仅约为一次矩阵向量乘法. 这一实现是 Eisenstat 技巧之一 \BookCite{53}.}
% block: ssor-preconditioner page:173 type:ordinary-paragraph
{对称逐次超松弛 (SSOR) 预条件子可视为 SGS 预条件子的推广. 该预条件子定义为}
\[K_{SSOR}:=(\widetilde D-L)\widetilde D^{-1}(\widetilde D-U),\quad\widetilde D=D/\omega.\]
{注意, 选取 $\omega=1$ 就得到 $K_{SGS}$. 此时 Eisenstat 技巧如下: 由 $A=D-L-U$ 和 $\widetilde A=(\widetilde D-L)^{-1}A(\widetilde D-U)^{-1}$, 有}
\[\begin{aligned}\widetilde A\bm z&=(\widetilde D-L)^{-1}A(\widetilde D-U)^{-1}\bm z\\
&=(\widetilde D-L)^{-1}[(\widetilde D-L)+(D-2\widetilde D)+(\widetilde D-U)](\widetilde D-U)^{-1}\bm z\\
&=(\widetilde D-U)^{-1}\bm z+(\widetilde D-L)^{-1}(D-2\widetilde D)(\widetilde D-U)^{-1}\bm z+(\widetilde D-L)^{-1}\bm z\\
&=\bm t+(\widetilde D-L)^{-1}[(D-2\widetilde D)\bm t+\bm z],\end{aligned}\]
{其中 $\bm t=(\widetilde D-U)^{-1}\bm z$.}
% block: preconditioner-reordering-heading pages:173,174 type:subsection
\subsubsection{{预条件子的重排序}}\label{subsec:preconditioner-reorderings}
{前面各小节考虑了用于有效求解原方程组的一些预条件技术. 另一方面, 自然也可以寻找有效的重排序以改善预条件子的性能, 例如, 对采用重排序矩阵 $P$ 的下列方程组应用 Krylov 子空间方法:}
\[K^{-1}PAP^{\top}\widetilde{\bm x}=\widetilde{\bm b},\]
{其中 $\widetilde{\bm x}=P\bm x$, $\widetilde{\bm b}=K^{-1}P\bm b$. 基于这一思想, 可以采用许多重排序技术, 如 Cuthill--Mckee (CM) \BookCite{42}, 反向 Cuthill--Mckee (RCM) \BookCite{75}, 嵌套剖分 (ND) \BookCite{76} 和最小度 (MD) \BookCite{77}. 关于这些算法, 另见 \BookCite{51}, \BookCite{123}, \BookCite{151}. CM 和 RCM 减小矩阵的带宽, MD 减少矩阵的填充元数, 而 ND 生成近似块对角矩阵, 适合并行计算.}
% block: reordering-evidence page:174 type:ordinary-paragraph
{对于对称定号矩阵, Duff 和 Meurant 给出了关于重排序效果的重要结果, 即用于直接求解器的重排序技术并未提高预条件子的质量 \BookCite{52}.}
{对于非对称矩阵, Saad 通过实验表明, 在 ILUT 预条件之前进行 MD 和 ND 重排序, 相比原方程组没有优势, 只有 RCM 最有可能带来改进 \BookCite{151}, p.334. 另一方面, Benzi 等 \BookCite{23} 研究了重排序对 ILU 型预条件子的影响, 得出在系数矩阵高度非对称的情形下 RCM 显著改善预条件子的结果. 文献 \BookCite{25} 也讨论了对近似逆预条件子的影响的类似结果.}
